\chapter{Dasar Teori}

\section{Konsep Multi-Stage Build dan Optimasi Citra Docker}

\subsection{Paradigma Multi-Stage Build dan Pemisahan Lingkungan}
Pada implementasi kontainerisasi konvensional (\textit{single-stage build}), seluruh perkakas pembangun (\textit{build tools}), manajer paket, dependensi pengembangan, serta kode sumber berada dalam satu citra yang sama. Pendekatan ini menghasilkan ukuran citra (\textit{image size}) yang sangat besar, memperlambat proses transfer melalui jaringan, serta memperluas bidang serangan (\textit{attack surface}) akibat banyaknya paket biner yang sebenarnya tidak dibutuhkan saat aplikasi berjalan [\ref{turnbull_docker}].

Untuk mengatasi keterbatasan tersebut, Docker memperkenalkan paradigma \textit{multi-stage build} [\ref{docker_multistage}]. Melalui mekanisme ini, pengembang dapat mendefinisikan beberapa instruksi \texttt{FROM} yang masing-masing merepresentasikan tahapan (\textit{stage}) mandiri dalam satu berkas \texttt{Dockerfile}. Setiap tahapan dapat menggunakan \textit{base image} yang berbeda sesuai fungsinya:
\begin{itemize}
    \item \textbf{Tahap Pembangun (\textit{Builder Stage}):} Dikonfigurasi dengan lingkungan lengkap yang memuat perkakas kompilasi berat (seperti manajer paket Composer, Node.js, atau SDK) untuk memproses dependensi dan mengompilasi artefak aplikasi.
    \item \textbf{Tahap Produksi (\textit{Runtime Stage}):} Menggunakan citra minimalis yang hanya menyalin artefak matang dari tahap pembangun melalui instruksi selektif \texttt{COPY --from=<stage>}, tanpa mengikutsertakan perkakas kompilasi asal.
\end{itemize}

\begin{figure}[H]
    \centering
    \includegraphics[width=0.65\textwidth]{multistage-architecture.png}
    \caption{Arsitektur konseptual multi-stage build: pemisahan artefak pembangun dari lingkungan produksi}
\end{figure}

\subsection{Strategi Reduksi Ukuran Citra dengan Base Image Minimalis}
Selain memisahkan tahap kompilasi, efisiensi citra akhir sangat ditentukan oleh pemilihan distribusi sistem operasi dasar (\textit{base image}). Distribusi standar berbasis Debian (seperti varian \texttt{php:8.3-cli}) menyertakan ribuan paket sistem berkas umum sehingga ukurannya membengkak hingga ratusan megabita [\ref{docker_multistage}]. Sebaliknya, varian berbasis Alpine Linux (seperti \texttt{php:8.3-cli-alpine}) mengadopsi pustaka sistem \textit{musl libc} dan antarmuka biner \textit{BusyBox}, menghasilkan fondasi sistem operasi yang sangat ringkas dengan ukuran dasar kurang dari 10 MB. Penggunaan Alpine secara empiris mampu memangkas volume citra akhir lebih dari 60\% sekaligus mengurangi potensi kerentanan keamanan bawaan (\textit{Common Vulnerabilities and Exposures}/CVE) sistem operasi [\ref{nist_container}].

\clearpage
\section{Aspek Keamanan dan Pemantauan Keandalan Kontainer}

\subsection{Prinsip Least Privilege dan Eksekusi Non-Root}
Secara bawaan (\textit{default}), proses di dalam kontainer Docker dijalankan dengan hak akses pengguna \textit{root} (\textit{UID 0}). Meskipun kontainer terisolasi di dalam \textit{namespace} Linux, eksekusi sebagai \textit{root} mengandung risiko eskalasi hak akses (\textit{privilege escalation}) dan kebocoran kontainer (\textit{container breakout}) apabila terjadi eksploitasi celah keamanan pada kernel sistem operasi \textit{host} [\ref{nist_container}].

Pedoman keamanan \textit{NIST Special Publication 800-190} merekomendasikan penerapan prinsip hak istimewa terendah (\textit{least privilege}) dengan menjalankan aplikasi kontainer di bawah pengguna non-\textit{root} [\ref{nist_container}]. Pada aplikasi web berbasis PHP, pengguna standar seperti \texttt{www-data} (\textit{UID 82} pada Alpine) didefinisikan melalui instruksi deklaratif \texttt{USER www-data}. Hak kepemilikan direktori kerja aplikasi dialihkan secara ketat (\texttt{chown -R www-data:www-data}) agar proses web memiliki wewenang membaca berkas sumber dan memanipulasi berkas lokal (seperti basis data SQLite dan direktori \textit{cache}) tanpa memiliki wewenang administratif terhadap sistem kontainer secara keseluruhan.

\begin{figure}[H]
    \centering
    \includegraphics[width=0.6\textwidth]{container-security.png}
    \caption{Model isolasi keamanan hak akses non-root (USER www-data) di dalam kontainer}
\end{figure}

\subsection{Mekanisme Health Monitoring melalui HEALTHCHECK}
Keandalan sistem perangkat lunak modern menuntut adanya pemantauan status kesehatan aplikasi secara berkesinambungan. Docker menyediakan instruksi \texttt{HEALTHCHECK} untuk menguji apakah layanan internal kontainer tidak hanya sekadar berjalan pada tabel proses, melainkan benar-benar mampu merespons permintaan transaksi data secara fungsional [\ref{docker_multistage}].

Pada ekosistem Laravel 11, telah disediakan rute pemeriksaan kesehatan terdedikasi (\texttt{/up}) yang secara bawaan memverifikasi kesiapan \textit{framework} dan dependensi krusial [\ref{laravel_health}]. Instruksi \texttt{HEALTHCHECK} mengonfigurasi probe berkala menggunakan parameter:
\begin{itemize}
    \item \texttt{--interval}: Jarak waktu antar-pengujian (misalnya setiap 15 detik).
    \item \texttt{--timeout}: Batas toleransi respons perintah probe (misalnya 3 detik).
    \item \texttt{--start-period}: Periode inisialisasi awal agar proses \textit{startup} kontainer tidak langsung dianggap gagal (misalnya 5 detik).
    \item \texttt{--retries}: Jumlah kegagalan berturut-turut sebelum status kontainer diubah menjadi \texttt{unhealthy} (misalnya 3 kali).
\end{itemize}
Mekanisme ini memungkinkan orkestrator kontainer atau sistem pemantauan untuk melakukan pemulihan otomatis (\textit{self-healing}) saat layanan mengalami degradasi.

\clearpage
\section{Otomatisasi CI/CD dan GitHub Container Registry (GHCR)}

\subsection{Arsitektur GitHub Actions untuk Publishing Citra Kontainer}
\textit{Continuous Integration and Continuous Deployment} (CI/CD) adalah praktik rekayasa perangkat lunak di mana setiap perubahan kode diuji, dikompilasi, dan dikirimkan secara otomatis menuju repositori artefak atau lingkungan produksi. GitHub Actions mengimplementasikan konsep ini melalui alur kerja berbasis peristiwa deklaratif (\textit{event-driven workflow}) yang dieksekusi oleh mesin pelaksana (\textit{runner}) terisolasi [\ref{ghcr_docs}].

\begin{figure}[H]
    \centering
    \includegraphics[width=0.7\textwidth]{ghcr-pipeline.png}
    \caption{Alur otomatisasi pipeline CI/CD GitHub Actions menuju GitHub Container Registry (GHCR)}
\end{figure}

GitHub Container Registry (GHCR) bertindak sebagai penyimpan paket citra OCI (\textit{Open Container Initiative}) yang terintegrasi secara mendalam dengan repositori kode GitHub [\ref{ghcr_docs}]. Proses otomatisasi publikasi citra memanfaatkan aksi resmi seperti \texttt{docker/login-action} untuk autentikasi nir-kata sandi menggunakan token bawaan \texttt{GITHUB\_TOKEN} serta \texttt{docker/build-push-action} untuk proses kompilasi citra multi-arsitektur. Penandaan versi (\textit{image tagging}) wajib mengadopsi nilai hash komit (\texttt{\$\{\{ github.sha \}\}}) guna menjamin keterlacakan (\textit{traceability}) dan kekekalan versi (\textit{immutability}) artefak rilis.

\subsection{Mekanisme Hak Akses dan Environment Protection Gates}
Publikasi paket ke GHCR memerlukan penegasan hak akses (\textit{scope permissions}) pada tingkat pekerjaan alur kerja, di mana izin \texttt{packages: write} wajib didefinisikan secara eksplisit [\ref{ghcr_docs}]. Tanpa izin ini, proses pengunggahan lapisan citra akan ditolak oleh sistem keamanan GitHub.

Selain itu, dalam siklus rilis berkelanjutan berstandar industri, penggelaran kode ke lingkungan produksi dilindungi oleh mekanisme gerbang perlindungan lingkungan (\textit{Environment Protection Rules}) [\ref{ghcr_docs}]. Aturan ini menahan (\textit{hold}) eksekusi alur kerja pada tahap penyebaran produksi hingga peninjau resmi (\textit{reviewer}) melakukan verifikasi manual dan memberikan persetujuan (\textit{approval}). Integrasi gerbang ini menjamin bahwa seluruh pengujian otomatis, kompilasi citra, dan peninjauan kualitas telah terpenuhi sepenuhnya sebelum aplikasi berdampak pada pengguna akhir.


\chapter{Hasil dan Pembahasan}

\section{Restrukturisasi Dockerfile ke Multi-Stage Build}

Pada praktikum sebelumnya (Tugas 4), proses pembangunan \textit{Docker image} menggunakan arsitektur \textit{single-stage build} berbasis \texttt{php:8.3-cli}. Pendekatan tersebut menghasilkan berkas citra berukuran besar karena seluruh perkakas kompilasi, dependensi sistem operasi Debian, dan binari Composer ikut terdistribusi ke dalam citra akhir. Pada praktikum ini, \texttt{Dockerfile} direstrukturisasi secara fundamental mengadopsi paradigma \textit{multi-stage build} guna memisahkan lingkungan kompilasi dependensi dari lingkungan eksekusi produksi.

\subsection{Definisi Stage Builder (composer:2.7)}
Tahap pertama didefinisikan sebagai \textit{stage} \texttt{builder} menggunakan \textit{base image} resmi \texttt{composer:2.7}. Tahap ini didedikasikan secara eksklusif untuk menyelesaikan resolusi paket dan instalasi dependensi vendor aplikasi Laravel. Strategi \textit{layer caching} diterapkan dengan menyalin berkas manifes dependensi \texttt{composer.json} dan \texttt{composer.lock} terlebih dahulu ke dalam direktori kerja \texttt{/app}. Eksekusi perintah \texttt{composer install} dijalankan dengan flag optimasi produksi: \texttt{--no-interaction}, \texttt{--prefer-dist}, \texttt{--ignore-platform-reqs}, \texttt{--no-scripts}, dan \texttt{--optimize-autoloader}. Pendekatan ini memastikan instalasi pustaka dependensi terisolasi penuh dan dapat di-\textit{cache} ulang secara efisien tanpa terpengaruh oleh modifikasi berkas kode program harian.

\begin{figure}[H]
    \centering
    \includegraphics[width=0.55\textwidth]{1.png}
    \caption{Konfigurasi tahap pertama (builder) pada Dockerfile berbasis image composer:2.7}
\end{figure}

\subsection{Konfigurasi Stage Runtime Berbasis Alpine, Non-Root User, dan HEALTHCHECK}
Tahap kedua merupakan citra produksi akhir yang dirancang seringkas dan seaman mungkin dengan menggunakan \textit{base image} minimalis berbasis Alpine Linux, yaitu \texttt{php:8.3-cli-alpine}. Citra ini hanya menginstalasi dependensi sistem esensial \texttt{sqlite-dev} melalui manajer paket \texttt{apk} tanpa menyimpan \textit{cache} lokal (\texttt{--no-cache}), serta mengompilasi ekstensi PHP minimal yang dibutuhkan aplikasi (\texttt{pdo\_mysql} dan \texttt{pdo\_sqlite}). 

Artefak yang telah matang dari tahap \textit{builder} disalin secara selektif menggunakan instruksi deklaratif \texttt{COPY --from=builder /app /var/www}. Demi memenuhi prinsip keamanan \textit{least privilege}, kepemilikan seluruh berkas dialihkan ke \textit{user} bawaan \texttt{www-data} melalui perintah \texttt{chown -R www-data:www-data /var/www}, dan kontainer diinstruksikan untuk beroperasi di bawah hak akses \texttt{USER www-data} (bukan \textit{root}). Selain itu, diimplementasikan mekanisme pemantauan keandalan menggunakan instruksi \texttt{HEALTHCHECK} yang secara periodik (setiap 15 detik) mengeksekusi \texttt{wget} ke rute kesehatan bawaan Laravel 11 (\texttt{http://127.0.0.1:8000/up}).

\begin{figure}[H]
    \centering
    \includegraphics[width=0.55\textwidth]{2.png}
    \caption{Konfigurasi tahap produksi pada Dockerfile dengan base image php:8.3-cli-alpine, USER www-data, dan instruksi HEALTHCHECK}
\end{figure}

\clearpage
\section{Eksekusi Build Citra dan Komparasi Reduksi Ukuran}

\subsection{Proses Inisiasi dan Kompilasi Multi-Stage Build}
Pembangunan citra baru diinisiasi melalui terminal dengan mengeksekusi perintah \texttt{docker build -t evolusi-api:tugas5 .}. Docker BuildKit secara cerdas menguraikan dependensi antar-\textit{stage} deklaratif. Pada tahap awal, BuildKit memicu pengunduhan manifes serta ekstraksi lapisan untuk kedua \textit{base image} yang dibutuhkan, yaitu \texttt{composer:2.7} dan \texttt{php:8.3-cli-alpine}.

\begin{figure}[H]
    \centering
    \includegraphics[width=0.8\textwidth]{5.png}
    \caption{Inisiasi eksekusi docker build tugas5 dan pengunduhan metadata base image}

    \vspace{0.5cm}

    \includegraphics[width=0.8\textwidth]{6.png}
    \caption{Tahap penyelesaian multi-stage build: kompilasi builder, transfer artefak vendor, dan perakitan citra akhir}
\end{figure}

Tahap pembangunan diselesaikan dalam total durasi 189,9 detik (17 langkah). Sebagaimana terekam pada log terminal, langkah \texttt{[builder 1/5]} hingga \texttt{[builder 5/5]} mengeksekusi instalasi dependensi vendor pada direktori sementara \texttt{/app}. Selanjutnya, langkah \texttt{[stage-1 4/6]} melakukan transfer artefak vendor matang ke \texttt{/var/www}, dilanjutkan pengaturan hak akses berkas \texttt{www-data}, pembangkitan kunci enkripsi aplikasi, dan migrasi skema basis data sebelum dilakukan penamaan manifes akhir ke \texttt{docker.io/library/evolusi-api:tugas5}.

\subsection{Analisis Reduksi Ukuran Citra Sebelum dan Sesudah Optimasi}
Verifikasi efektivitas optimasi dilakukan dengan membandingkan citra hasil tugas sebelumnya (\texttt{evolusi-api:tugas4}) dengan citra hasil restrukturisasi \textit{multi-stage} (\texttt{evolusi-api:tugas5}) menggunakan perintah \texttt{docker images} di terminal dan inspeksi visual pada panel Docker Desktop.

\begin{figure}[H]
    \centering
    \includegraphics[width=0.85\textwidth]{7.png}
    \caption{Verifikasi ukuran citra lokal melalui perintah docker images}
\end{figure}

Pemeriksaan komparatif juga dilakukan secara visual melalui panel \textit{Images} pada Docker Desktop untuk mengonfirmasi status keaktifan kontainer dan ukuran fisik citra secara berdampingan.

\begin{figure}[H]
    \centering
    \includegraphics[width=0.9\textwidth]{17.png}
    \caption{Komparasi ukuran citra tugas4 dan tugas5 pada antarmuka Docker Desktop Images}
\end{figure}

Berdasarkan data observasi empiris pada terminal dan Docker Desktop, diperoleh perbandingan ukuran citra sebagai berikut:
\begin{enumerate}
    \item \textbf{Ukuran Penggunaan Disk (\textit{Disk Usage}):}
    \begin{itemize}
        \item Sebelum optimasi (\texttt{evolusi-api:tugas4}): 949,94 MB (dibulatkan menjadi 950 MB).
        \item Sesudah optimasi (\texttt{evolusi-api:tugas5}): 336,72 MB (dibulatkan menjadi 337 MB).
        \item \textbf{Selisih Penurunan Ukuran:} $949{,}94\text{ MB} - 336{,}72\text{ MB} = \mathbf{613{,}22\text{ MB}}$.
        \item \textbf{Persentase Reduksi:}
        $$\frac{613{,}22\text{ MB}}{949{,}94\text{ MB}} \times 100\% = \mathbf{64{,}55\%}$$
    \end{itemize}
    \item \textbf{Ukuran Konten Lapisan (\textit{Content Size}):}
    \begin{itemize}
        \item Sebelum optimasi (\texttt{tugas4}): 234 MB.
        \item Sesudah optimasi (\texttt{tugas5}): 64,9 MB.
        \item \textbf{Selisih Penurunan Konten:} $234\text{ MB} - 64{,}9\text{ MB} = \mathbf{169{,}1\text{ MB}}$ (reduksi sebesar $\mathbf{72{,}26\%}$).
    \end{itemize}
\end{enumerate}
Penurunan drastis sebesar 64,55\% pada penyimpanan disk dan 72,26\% pada ukuran konten membuktikan bahwa transisi ke \textit{base image} Alpine serta eliminasi dependensi pembangun Composer dari citra akhir sangat efektif dalam memangkas ukuran distribusi kontainer.

\clearpage
\subsection{Monitoring Durasi dan Caching Layer pada Docker Desktop BuildKit}
Menu \textit{Builds} pada Docker Desktop memberikan visualisasi komparatif terhadap histori sesi pembangunan kontainer. Sesi pembangunan \textit{multi-stage build} pertama (\textit{cold build}) dengan ID \texttt{pwdfdd} tercatat memakan durasi 3 menit 9 detik (189,9 detik), disebabkan oleh waktu unduh binari paket kompilasi Alpine dan resolusi Composer. Riwayat ini tersimpan berdampingan dengan sesi rilis Tugas 4 sebelumnya (durasi \textit{build} 1 menit 27 detik hingga pengujian \textit{cache hit} 1,4 detik).

\begin{figure}[H]
    \centering
    \includegraphics[width=0.9\textwidth]{16.png}
    \caption{Histori pencatatan sesi build pada antarmuka Docker Desktop Builds}
\end{figure}

\clearpage
\section{Verifikasi Runtime, Keamanan, dan Fungsionalitas Kontainer}

\subsection{Pengujian Hak Akses Non-Root dan Status Keaktifan Kontainer}
Untuk membuktikan kepatuhan terhadap batasan keamanan dan operasional, kontainer diinstansiasi di latar belakang (\textit{detached mode}) menggunakan perintah \texttt{docker run -d --name laravel-api -p 8000:8000 evolusi-api:tugas5}. Kontainer menghasilkan ID hash \texttt{25f99bda308c...}.

Pengujian identitas pengguna eksekusi dilakukan secara langsung ke dalam kontainer melalui perintah \texttt{docker exec laravel-api whoami}. Keluaran terminal menampilkan nilai \texttt{www-data}, yang membuktikan secara valid bahwa proses utama kontainer tidak berjalan di bawah hak istimewa \textit{root}. Selanjutnya, pemeriksaan status operasional kontainer melalui perintah \texttt{docker ps} memperlihatkan kontainer aktif dengan status \texttt{Up 51 seconds (healthy)}. Status \texttt{(healthy)} tersebut mengonfirmasi bahwa instruksi \texttt{HEALTHCHECK} berhasil menghubungi rute internal \texttt{/up} dan mengembalikan kode status HTTP sukses secara berulang.

\begin{figure}[H]
    \centering
    \includegraphics[width=0.9\textwidth]{8.png}
    \caption{Verifikasi eksekusi user non-root (www-data) via whoami dan status operasional healthy pada docker ps}
\end{figure}

\clearpage
\section{Konfigurasi dan Otomatisasi Workflow GHCR}

\subsection{Perancangan Workflow GitHub Actions untuk GHCR}
Penyusunan berkas alur kerja otomatisasi diwujudkan melalui berkas \texttt{.github/workflows/\allowbreak docker-publish.yml}. Alur kerja ini dirancang untuk terpicu secara otomatis pada setiap aksi \textit{push} menuju \textit{branch} \texttt{main} dan \texttt{dev}. Konfigurasi penting pada level pekerjaan (\textit{job}) \texttt{build-and-push} mencakup penetapan blok perizinan (\textit{permissions}) secara eksplisit, yaitu \texttt{contents: read} untuk membaca kode sumber serta \texttt{packages: write} sebagai syarat mutlak agar GitHub Actions memiliki wewenang mengunggah citra ke GitHub Container Registry (GHCR).

\begin{figure}[H]
    \centering
    \includegraphics[width=0.7\textwidth]{3.png}
    \caption{Konfigurasi workflow docker-publish.yml mendefinisikan pemicu event dan perizinan packages: write}
\end{figure}

\subsection{Troubleshooting dan Penanganan Isu Scope Environment Variable}
Pada implementasi tahapan pembangunan dan publikasi citra, terdapat aturan penamaan GHCR yang mewajibkan nama repositori/citra ditulis dalam format huruf kecil penuh (\textit{lowercase}). Langkah \texttt{Lowercase Repository Name} dikonfigurasi untuk mengonversi \texttt{\$\{\{ github.repository \}\}} menjadi huruf kecil menggunakan utilitas shell \texttt{tr '[:upper:]' '[:lower:]'} dan menyimpannya ke berkas \texttt{\$GITHUB\_ENV}. 

Namun, pada saat dilakukan pemanggilan variabel lingkungan \texttt{\$\{\{ env.IMAGE\_NAME \}\}} di dalam blok argumen \texttt{tags} pada aksi \texttt{docker/build-push-action@v5}, sistem analisis linter memberikan peringatan (\textit{warning}) bergaris bawah kuning bertuliskan \textit{``Context access might be invalid: env.IMAGE\_NAME''}. Peringatan ini muncul karena variabel yang diinjeksikan secara dinamis melalui \texttt{\$GITHUB\_ENV} pada \textit{runtime} suatu langkah baru dievaluasi pada saat eksekusi langkah berikutnya, sehingga tidak terdeteksi pada fase validasi skema statis. Penanganan dilakukan dengan memastikan urutan langkah beroperasi secara sekuensial dan variabel dievaluasi dengan tepat pada \textit{runtime step}. Selain itu, parameter \texttt{tags} disetel untuk menggunakan hash unik commit (\texttt{\$\{\{ github.sha \}\}}) sebagai penanda versi (\textit{tag}) yang kekal (\textit{immutable}), bukan sekadar menggunakan tag generik \texttt{latest}.

\begin{figure}[H]
    \centering
    \includegraphics[width=0.65\textwidth]{4.png}
    \caption{Peringatan konteks env.IMAGE\_NAME pada konfigurasi tags aksi docker/build-push-action}
\end{figure}

\subsection{Eksekusi Pipeline CI/CD dan Publikasi Citra pada Branch Fitur}
Pengujian integrasi alur kerja dimulai saat pengembang melakukan \textit{push} pembaruan kode pada \textit{branch} fitur \texttt{feature/docker-multistage}. Di repositori organisasi pribadi (\texttt{AcademyCity-RND}), \textit{pipeline} otomatis mengeksekusi rangkaian \textit{check}: \texttt{Build Application} serta \texttt{Automated Unit \& Feature Tests}. Seluruh pengujian berhasil dilalui dalam waktu 50 detik dengan status centang hijau (\textit{Success}). Sementara itu, repositori organisasi kelas (\texttt{KEPL2026}) menerima kode sumber yang identik dan menampilkan notifikasi kesiapan pembuatan \textit{Pull Request}.

\begin{figure}[H]
    \centering
    \includegraphics[width=0.9\textwidth]{9.png}
    \caption{Eksekusi pipeline CI berhasil pada branch fitur (Kiri: repo pribadi AcademyCity-RND, Kanan: repo KEPL2026)}
\end{figure}

\clearpage
\section{Manajemen Repositori, Alur Pull Request, dan Deployment Gating}

Manajemen repositori pada praktikum ini menerapkan arsitektur kolaborasi bercabang dengan dualitas lingkungan organisasi. Berdasarkan status resmi organisasi mata kuliah (\texttt{KEPL2026}) yang sedang dalam peninjauan aju banding sehingga menonaktifkan fitur GitHub Actions, seluruh otomasi pengujian, kompilasi Docker, dan publikasi GHCR dieksekusi melalui repositori organisasi pribadi (\texttt{AcademyCity-RND}), sementara rekam jejak kode tetap dipelihara dan disinkronkan secara konsisten di kedua repositori.

\subsection{Alur Pull Request Feature ke Dev dan Sinkronisasi Dual Repository}
Promosi kode dari \textit{branch} pengembangan fitur menuju \textit{branch} integrasi \texttt{dev} diajukan melalui \textit{Pull Request} \#15 bertajuk \textit{``feat: Multi-stage Docker, Alpine, Non-root \& GHCR (Tugas 5)''}. Pada repositori pribadi, PR \#15 mencatat 6 commit dengan validasi otomatis sebanyak 13 \textit{checks} CI/CD berstatus lulus. Pada repositori KEPL2026, PR mencakup 4 commit tanpa eksekusi \textit{checks} akibat restriksi organisasi, namun berstatus bersih (\textit{``No conflicts with base branch''}) dan siap digabungkan.

\begin{figure}[H]
    \centering
    \includegraphics[width=0.9\textwidth]{10.png}
    \caption{Pull Request \#15 integrasi fitur Tugas 5 dari branch feature ke dev (Kiri: repo pribadi, Kanan: repo KEPL2026)}
\end{figure}

\clearpage
\subsection{Promosi Rilis Menuju Branch Main (Pull Request \#16)}
Tahap rilis final dilakukan dengan membuat \textit{Pull Request} \#16 bertajuk \textit{``Release: Tugas 5 (Multi-stage Docker \& GHCR)''} dari \textit{branch} \texttt{dev} menuju \texttt{main}. Penggabungan ini merangkum seluruh perbaikan konfigurasi, sinkronisasi skema basis data SQLite, penyesuaian dependensi Alpine, pembaruan dokumentasi README, serta integrasi alur kerja rilis GHCR.

\begin{figure}[H]
    \centering
    \includegraphics[width=0.9\textwidth]{11.png}
    \caption{Pull Request \#16 rilis Tugas 5 menuju branch main (Kiri: repo pribadi dengan 16 commit, Kanan: repo KEPL2026)}
\end{figure}

\subsection{Validasi Pipeline CI/CD dan Deployment Environment Protection}
Saat PR \#16 digabungkan ke \textit{branch} \texttt{main}, GitHub Actions secara otomatis memicu seluruh alur kerja integrasi berkelanjutan. Tinjauan status \textit{checks} menunjukkan bahwa 8 pekerjaan (\textit{jobs}) berhasil diselesaikan dengan sempurna, termasuk pengujian unit, kompilasi aplikasi, \textit{linter}, \textit{deployment} lingkungan \textit{staging}, serta pekerjaan \texttt{Publish Docker Image to GHCR / build-and-push} yang tuntas dalam durasi 46 detik. 

Satu pekerjaan tersisa, yaitu \texttt{Deploy to Production}, ditahan dalam status menunggu (\textit{Waiting}). Hal ini merupakan implementasi nyata dari fitur \textit{Environment Protection Rules} pada GitHub Enterprise/Actions, di mana proses rilis ke lingkungan produksi membutuhkan persetujuan manual (\textit{manual approval}) dari \textit{reviewer} yang ditunjuk sebelum instruksi penggelaran dieksekusi.

\begin{figure}[H]
    \centering
    \includegraphics[width=0.9\textwidth]{12.png}
    \caption{Tampilan ringkasan checks pada branch main: 8 checks sukses dan 1 check rilis produksi menunggu persetujuan}

    \vspace{0.5cm}

    \includegraphics[width=0.9\textwidth]{13.png}
    \caption{Modal antarmuka Review pending deployments untuk otorisasi manual rilis ke lingkungan production}
\end{figure}

Setelah \textit{reviewer} memberikan persetujuan melalui tombol \textit{``Approve and deploy''}, pekerjaan rilis produksi segera dijalankan hingga selesai. Hasil sinkronisasi akhir pada kedua repositori memperlihatkan status commit terbaru \texttt{Merge pull request \#16} pada \textit{branch} \texttt{main}, dengan lencana \textit{Deployments: production} berstatus aktif hijau pada repositori pribadi.

\begin{figure}[H]
    \centering
    \includegraphics[width=0.9\textwidth]{14.png}
    \caption{Kondisi akhir branch main pada kedua repositori setelah proses rilis Tugas 5 sukses terintegrasi}
\end{figure}

\subsection{Verifikasi Manajemen Paket dan Visibilitas Publik pada GHCR}
Hasil publikasi kontainer diverifikasi secara langsung pada halaman manajemen paket GitHub Container Registry organisasi \texttt{AcademyCity-RND}. Paket dengan nama \texttt{evolusi-pl-24-542159-sv-25009-v2} berhasil terbit dan menampilkan instruksi unduhan berbasis baris perintah:
\begin{center}
\texttt{docker pull ghcr.io/academycity-rnd/\allowbreak evolusi-pl-24-542159-sv-25009-v2:\allowbreak ddc450a5143d842c62bdf70693dcf1ed08512c7b}
\end{center}
Tag citra yang diterbitkan secara presisi merujuk pada hash unik commit \texttt{ddc450a...}. Sesuai dengan instruksi penugasan, visibilitas paket telah diubah dan dipastikan berstatus publik (\textit{Public}) sehingga dapat diakses dan diunduh secara bebas oleh publik maupun sistem orkestrasi eksternal.

\begin{figure}[H]
    \centering
    \includegraphics[width=0.9\textwidth]{15.png}
    \caption{Halaman rilis package pada GitHub Container Registry (GHCR) menunjukkan tag berbasis commit hash dan status visibilitas publik}
\end{figure}

\section{Tautan Repositori Kode Sumber dan Registry Package}

Seluruh berkas konfigurasi, kode program aplikasi, dan artefak kontainer praktikum ini tersedia secara terbuka pada tautan berikut:
\begin{itemize}
    \item Repositori Matkul KEPL2026 (v1):\\
    \url{https://github.com/KEPL2026/evolusi-pl-24-542159-SV-25009-v1}

    \item Repositori Pribadi AcademyCity-RND (v2 - GitHub Actions Active):\\
    \url{https://github.com/AcademyCity-RND/evolusi-pl-24-542159-SV-25009-v2}

    \item Halaman Paket GitHub Container Registry (GHCR):\\
    \url{https://github.com/AcademyCity-RND/evolusi-pl-24-542159-SV-25009-v2/pkgs/container/evolusi-pl-24-542159-sv-25009-v2}
\end{itemize}


\chapter{Kesimpulan}

Berdasarkan praktikum mengenai optimasi citra Docker melalui pendekatan \textit{multi-stage build}, peningkatan aspek keamanan kontainer, dan otomatisasi alur kerja CI/CD menuju GitHub Container Registry (GHCR) yang telah dilakukan, dapat ditarik beberapa kesimpulan utama sebagai berikut:
\begin{enumerate}
    \item Restrukturisasi \texttt{Dockerfile} dari pendekatan \textit{single-stage build} menjadi \textit{multi-stage build} berhasil memisahkan tahap kompilasi dependensi (berbasis \texttt{composer:2.7}) dari tahap eksekusi produksi (berbasis \texttt{php:8.3-cli-alpine}), sehingga perkakas berat Composer dan dependensi pembangun tidak terdistribusi ke dalam citra akhir.
    \item Penerapan \textit{multi-stage build} dan transisi ke \textit{base image} Alpine Linux secara empiris menghasilkan penurunan ukuran citra yang sangat signifikan:
    \begin{itemize}
        \item Penggunaan ruang penyimpanan disk (\textit{Disk Usage}) terpangkas dari 949,94 MB (\texttt{evolusi-api:tugas4}) menjadi 336,72 MB (\texttt{evolusi-api:tugas5}), menghasilkan selisih penghematan sebesar 613,22 MB atau efisiensi sebesar \textbf{64{,}55\%}.
        \item Ukuran konten lapisan (\textit{Content Size}) mengalami penurunan drastis dari 234 MB menjadi 64,9 MB, menghasilkan reduksi konten sebesar 169,1 MB atau efisiensi sebesar \textbf{72{,}26\%}.
    \end{itemize}
    \item Implementasi prinsip keamanan \textit{least privilege} berhasil dipenuhi dengan mendefinisikan direktif \texttt{USER www-data} dan pemindahan hak milik direktori kerja melalui perintah \texttt{chown -R www-data:www-data}, yang diverifikasi secara valid melalui keluaran perintah \texttt{docker exec laravel-api whoami} yang menampilkan identitas pengguna non-\textit{root} (\texttt{www-data}).
    \item Mekanisme pemantauan keandalan mandiri melalui instruksi \texttt{HEALTHCHECK} yang memantau rute kesehatan bawaan Laravel 11 (\texttt{http://127.0.0.1:8000/up}) setiap 15 detik terbukti beroperasi optimal, ditandai dengan status operasional kontainer yang berstatus \texttt{Up (healthy)} pada perintah \texttt{docker ps}.
    \item Otomatisasi rilis melalui GitHub Actions (\texttt{docker-publish.yml}) dengan deklarasi izin \texttt{packages: write} dan mekanisme penamaan huruf kecil (\textit{lowercase}) berhasil menerbitkan citra ke GitHub Container Registry (GHCR) secara publik dengan skema penandaan versi berbasis hash komit (\texttt{\$\{\{ github.sha \}\}}), serta dilindungi oleh mekanisme \textit{Environment Protection} rilis produksi yang mensyaratkan persetujuan manual (\textit{manual approval}) dari peninjau.
\end{enumerate}


\chapter{Daftar Pustaka}

\begin{enumerate}
    \item \label{turnbull_docker} J. Turnbull. \textit{The Docker Book: Containerization Is the New Virtualization}. James Turnbull, 2014. \mbox{[Online]. Tersedia:} \url{https://www.dockerbook.com/}

    \item \label{docker_multistage} Docker, Inc. \textit{Multi-stage builds: Use multi-stage builds}. Docker Documentation, 2024. \mbox{[Online]. Tersedia:} \url{https://docs.docker.com/build/building/multi-stage/}

    \item \label{nist_container} M. Souppaya, J. Morello, \& K. Scarfone. \textit{Application Container Security Guide}. NIST Special Publication 800-190, National Institute of Standards and Technology, 2017. \mbox{[Online]. Tersedia:} \url{https://doi.org/10.6028/NIST.SP.800-190}

    \item \label{laravel_health} Laravel LLC. \textit{Laravel 11.x Documentation: Routing and Application Health}. Laravel, 2024. \mbox{[Online]. Tersedia:} \url{https://laravel.com/docs/11.x/routing#routing-the-health-check-endpoint}

    \item \label{ghcr_docs} GitHub, Inc. \textit{Working with the Container Registry (GHCR)}. GitHub Docs, 2024. \mbox{[Online]. Tersedia:} \url{https://docs.github.com/en/packages/working-with-a-github-packages-registry/working-with-the-container-registry}

    \item \label{sourcecode} Tautan Repositori Kode Sumber dan Paket Publik:
    \begin{itemize}
        \item Repositori Matkul KEPL2026 (v1):\\
        \url{https://github.com/KEPL2026/evolusi-pl-24-542159-SV-25009-v1}

        \item Repositori Pribadi AcademyCity-RND (v2 - GitHub Actions Active):\\
        \url{https://github.com/AcademyCity-RND/evolusi-pl-24-542159-SV-25009-v2}

        \item Paket GitHub Container Registry (GHCR):\\
        \url{https://github.com/AcademyCity-RND/evolusi-pl-24-542159-SV-25009-v2/pkgs/container/evolusi-pl-24-542159-sv-25009-v2}
    \end{itemize}
\end{enumerate}
